%===============================================================================
% SEÇÃO 2 — Controle de Processos Endotérmicos com Reatores Tubulares
%===============================================================================

\section{Controle de Processos Endotérmicos}
\label{sec:controle_endotermico}

% -----------------------------------------------------------------------
% 2.1 — FORMULAÇÃO DO PROBLEMA DE CONTROLE
% -----------------------------------------------------------------------
\subsection{Formulação do problema de controle}
\label{sec:formulacao}

% TODO: verificar se existe figura genérica de reator tubular aquecido externamente
% em alguma referência clássica (ex.: Luyben 2014) para incluir nesta subsecção.
%
% *** ESPAÇO RESERVADO PARA FIGURA DO REATOR ***
% Substituir o bloco abaixo pela figura real quando disponível.
%
% \begin{figure}[htb]
%   \centering
%   \includegraphics[width=0.9\columnwidth]{figures/reator_esquematico}
%   \caption{Representação esquemática do reator tubular da planta B2H2/UFPR.
%            Indicar: zona de aquecimento (forno), leito catalítico, pontos
%            de medida de temperatura (TC) e posição do sensor de referência
%            para a malha de controle cascata.}
%   \label{fig:reator}
% \end{figure}

O processo de reforma a seco de biogás envolve um reator tubular de leito
fixo inserido em um forno elétrico, cuja dinâmica é descrita, de forma
genérica, pela equação de estado
\begin{equation}
  \dot{x}(t) = f\bigl(x(t),\,u(t),\,d(t)\bigr), \qquad y(t) = g\bigl(x(t)\bigr),
  \label{eq:estadogenerico}
\end{equation}
onde $x = [T_r,\; X_{\mathrm{CH}_4},\; r_{\mathrm{H}_2/\mathrm{CO}}]^\top$
é o vetor de estados (temperatura do leito, conversão de metano e razão
H\(_2\)/CO), $u = [P_{\mathrm{forno}},\; F_{\mathrm{CH}_4},\; F_{\mathrm{CO}_2}]^\top$
é o vetor de entradas manipuladas e $d$ representa perturbações
exernas (composição do biogás, temperatura ambiente,
pressão de operação).\citep{oliveira2021,luyben2014}

O problema de controle consiste em determinar $u(t)$ tal que o sistema
(\ref{eq:estadogenerico}) satisfaça, em regime permanente e durante transientes,
as seguintes restrições operacionais:\citep{luyben2014,thermodynamics2025}
\begin{align}
  T_r &\in [650,\;800]~^\circ\text{C}, \label{eq:restr_T} \\
  X_{\mathrm{CH}_4} &> 0{,}90, \label{eq:restr_X} \\
  1 &\leq F_{\mathrm{CO}_2}/F_{\mathrm{CH}_4} \leq 1{,}5. \label{eq:restr_razao}
\end{align}

O critério de desempenho adotado é o Índice de Erro Quadrático Integrado
ponderado, com penalização sobre os incrementos de controle:
\begin{equation}
  J = \int_0^\infty \Bigl[\|e(t)\|_Q^2 + \|\Delta u(t)\|_R^2 \Bigr]\,dt,
  \label{eq:ISE}
\end{equation}
onde $Q \succ 0$ e $R \succ 0$ são matrizes de ponderamento que refletem
relativa importância do erro de rastreamento frente ao custo de atuação.

% -----------------------------------------------------------------------
% 2.2 — DESAFIOS DO CONTROLE EM FORNOS COM REATORES TUBULARES
% -----------------------------------------------------------------------
\subsection{Desafios dinâmicos do forno-reator tubular}
\label{sec:desafios}

Reatores tubulares de leito fixo inseridos em fornos elétricos apresentam
características dinâmicas que dificultam significativamente o projeto de
malhas de controle convencionais.\citep{luyben2014} O primeiro intempérie é a
elevada inércia térmica do sistema: a massa de material refratário do forno,
alinhada com a massa do leito catalítico e das paredes do reator, resulta em
constantes de tempo térmicas da ordem de dezenas de minutos, tornando a
resposta lenta e prejudicando a rejeitação de perturbações de carga.\citep{luyben2014}

Alguns dos principais desafios operacionais desse tipo de configuração são:
\begin{itemize}
  \item \textbf{Atraso de transporte e distribuição de temperatura:} o perfil
    axial de temperatura ao longo do leito pode apresentar gradientes
    consideráveis, de modo que a temperatura medida por um único sensor não
    representa necessariamente o estado térmico global do leito.\citep{luyben2014}
  \item \textbf{Acoplamento térmico forno-reator:} a potência elétrica entregue
    ao forno afeta a temperatura da mufla, que por sua vez transfere calor
    ao reator por radiação e condução; esse acoplamento é altamente não
    linear e varia com a temperatura.\citep{luyben2014}
  \item \textbf{Não linearidade da reação:} a cinética da reforma a seco é
    fortemente não linear em temperatura, e a entalpia de reação
    (\(\approx 247~\text{kJ}\cdot\text{mol}^{-1}\)) atua como uma pertubação
    endotérmica de magnitude variável sobre o balanço de energia do forno,
    dificultando a manutenção de um ponto de operação fixo.\citep{thermodynamics2025}
  \item \textbf{Variabilidade da alimentação:} no caso do biogás, a
    composição CH\(_4\)/CO\(_2\) e a presença de inertes (N\(_2\), H\(_2\)S)
    variam ao longo do tempo, alterando o calor absorvido pela reação e
    deslocando o ponto de operação sem aviso prévio ao controlador.
    \citep{lavoie2014,thermodynamics2025}
  \item \textbf{Risco de formação de coque:} quedas de temperatura abaixo de
    640~\(^\circ\)C, ainda que transitórias, favorecem a deposição de
    carbono sólido, impondo assimetria nas consequências do erro de
    controle: o custo de um erro negativo (subaquecimento) é muito maior
    do que o de um erro positivo de mesma magnitude.\citep{thermodynamics2025,lavoie2014}
\end{itemize}

% -----------------------------------------------------------------------
% 2.3 — ESTRATÉGIA DE CONTROLE ADOTADA E SUA JUSTIFICATIVA
% -----------------------------------------------------------------------
\subsection{Estratégia de controle adotada}
\label{sec:estrategia}

Uma abordagem consagrada em aplicações industriais de fornos com reatores
tubulares é o uso de estruturas em cascata, nas quais uma malha interna
regula diretamente a variável associada ao fornecimento de calor enquanto
uma malha externa atua sobre a temperatura de saída do reator ou sobre um
ponto representativo do leito.\citep{luyben2014} Essa configuração melhora a
rejeição a perturbações de carga e de composição e reduz o efeito dos
atrasos associados à inércia térmica do conjunto forno-reator, sendo
precisamente a arquitetura adotada na planta B2H2/UFPR,
conforme detalhado na Seção~\ref{sec:processo}.

A literatura registra propostas de estratégias de maior complexidade
computacional para processos de reforma a seco, entre elas controle
preditivo baseado em modelos (MPC), estruturas neuro-fuzzy e
ANFIS.\citep{essien2019} Essas abordagens possuem mérito documentado em
plataformas computacionais adequadas, pois permitem lidar de forma explícita
com não linearidades e restrições operacionais. Entretanto, sua implementação
exige controladores lógicos programáveis (CLPs) ou computadores industriais
com capacidade de cómputo e memória significativamente superior àquela dos
equipamentos de custo acessível disponíveis para a plataforma experimental
B2H2/UFPR --- o que, no estágio atual de comissionamento da planta, inviabiliza
essa rota.\citep{essien2019,luyben2014}

Diante dessas restrições de hardware, a estratégia PID aprimorada --- cujas
modificações estruturais em relação ao controlador PID clássico constituem
o objeto central deste trabalho --- configura-se como a alternativa de
maior viabilidade prática para o nível de infraestrutura disponível,
assegurando comportamento previsível, confiável e de fácil manutenção
em campo. A eventual incorporação de estratégias avançadas permanece como
perspectiva de trabalho futuro, condicionada à disponibilidade de
hardware compatível e de modelos identificados em condições reais de
operação com biogás.\citep{essien2019}

% -----------------------------------------------------------------------
% 2.4 — ANÁLISE TERMODINÂMICA COMO SUPORTE AO CONTROLE
% -----------------------------------------------------------------------
\subsection{Análise termodinâmica como suporte ao controle}
\label{sec:revisao}

Análises termodinâmicas de equilíbrio têm sido amplamente empregadas para
mapear as janelas operacionais da DRM e da BDR em função de temperatura,
presão e razão CO\(_2\)/CH\(_4\).\citep{thermodynamics2025,lavoie2014} Esses
estudos identificam as regiões em que a conversão de CH\(_4\) e CO\(_2\) é
elevada e a formação de carbono sólido é termodinamicamente desfavorecida,
fornecendo limites superiores e inferiores diretos para as variáveis de
operação.\citep{thermodynamics2025} Na prática, esses limites se traduzem
em restrições concretas sobre os \emph{setpoints} das malhas de controle:
a temperatura do leito deve ser mantida acima de 640~\(^\circ\)C para garantir
espontaneidade reacional e abaixo de 800~\(^\circ\)C para preservar o
catalisador, enquanto a razão CO\(_2\)/CH\(_4\) deve permanecer próxima ou
acima da estequiometria unitária para desfavorecer a deposição de
coque.\citep{thermodynamics2025,lavoie2014} Essas informações fundamentam
tanto o projeto das malhas regulatórias quanto o desenvolvimento de esquemas
de proteção e, em níveis hierárquicos mais elevados, problemas de
otimização em tempo real.\citep{thermodynamics2025}
